\chapter{文献综述}

\section{负荷与净负荷预测方法}

电力负荷预测是电力系统运行与规划的基础环节，其方法论发展可划分为三个阶段。

\textbf{统计方法阶段。}以自回归积分滑动平均（ARIMA）和指数平滑为代表的时间序列模型长期占据负荷预测的主导地位。这类方法基于平稳性假设，通过自相关结构捕捉负荷的周期性和趋势性。Hong和Fan\cite{hong2016probabilistic}在对概率负荷预测的教程综述中指出，统计方法在日前和周前尺度上仍具竞争力，但在处理强非线性和多变量交互方面存在固有局限。

\textbf{机器学习阶段。}随机森林（Random Forest）、支持向量机（SVM）和梯度提升树（XGBoost）等方法引入了非线性建模能力和特征交互能力。这些方法的优势在于对异构特征（气象变量、日历变量、历史负荷）的灵活整合，且不要求输入数据满足特定分布假设\cite{wang2023hybrid}。然而，机器学习方法的预测精度依赖于人工特征工程的质量，且在超长序列依赖关系的建模上不如后续的深度学习方法。

\textbf{深度学习阶段。}LSTM通过门控机制解决了传统RNN的梯度消失问题，能够捕捉数百步的时间依赖关系。Silva-Rodriguez等\cite{silva2024lstm}将LSTM应用于含风光微电网的净负荷预测，在NAPS 2024上报告了优于传统方法的结果。Transformer架构通过多头自注意力机制建模全局时间模式，在长序列预测上展现出优势\cite{ebrahimzadeh2024ai}。CNN及其与RNN的混合架构——如时间卷积网络（TCN）——通过多尺度卷积核提取不同时间粒度的特征模式\cite{wang2023hybrid}。Ebrahimzadeh等\cite{ebrahimzadeh2024ai}在2024年对短期负荷预测AI技术的综述中系统比较了上述方法，指出深度学习在精度上的优势已被充分验证，但可解释性的缺失仍是制约其在调度实践中推广的主要障碍。

\textbf{净负荷预测的特殊挑战。}区别于传统负荷预测，净负荷（net load = load $-$ wind $-$ solar）的预测需要同时应对三个相互耦合的不确定性来源。Kaur等\cite{kaur2016net}是较早系统研究净负荷预测的工作，明确提出高可再生能源渗透率电网的净负荷预测问题，并指出风光出力的不确定性在净负荷上的叠加效应远超简单加总。Hanifi等\cite{hanifi2020critical}在对风电预测方法的综述中详细分析了风速的非平稳性和多尺度波动特征。Inman等\cite{inman2013solar}在太阳能预测综述中指出，辐照度受云量、气溶胶和太阳高度角的共同影响，其日内变异性在短时预测尺度上构成主要挑战。净负荷同时继承了这两种波动模式，加之负荷自身的复杂周期性（日周期、周周期、年周期及假日效应），使得净负荷预测成为一个在输入维度和物理机制上都更为复杂的任务。

综合来看，深度学习方法在净负荷预测的精度维度上已取得显著进展，但所有高精度方法均属黑箱模型，无法输出可审计的数学公式。如何在保持预测能力的同时实现物理层面的可解释性，是该领域的开放问题。

\section{符号回归与方程发现}

符号回归（symbolic regression, SR）旨在从数据中发现闭式数学表达式，与数值回归（如神经网络拟合）的根本区别在于：其输出不是参数化的黑箱函数，而是由基本数学算子（加、乘、三角函数、幂函数等）组成的、人类可读可验证的解析表达式。

\textbf{基于遗传编程的方法。}Schmidt和Lipson\cite{schmidt2009distilling}在2009年发表于Science的工作开创了现代数据驱动方程发现的先河。他们的Eureqa系统通过遗传编程（genetic programming, GP）在表达式树空间中搜索，成功从物理实验数据中重发现了牛顿运动定律和守恒律。这一工作确立了符号回归的基本范式：在由数学算子构成的离散组合空间中，以拟合精度和表达式复杂度为双目标，通过进化搜索找到Pareto前沿上的最优表达式。

后续发展中，PySR\cite{cranmer2023pysr}成为最广泛使用的开源符号回归工具。PySR采用多种群进化策略，结合常数优化和表达式简化，在SRBench基准测试\cite{lacava2021srbench}中展现了稳健的性能。Operon\cite{burlacu2020operon}以C++实现了高效的GP引擎，在速度和大规模数据处理上具有优势。Glyph\cite{balandat2020glyph}则面向工业应用场景提供了符号回归工具链。

\textbf{基于神经网络的方法。}AI Feynman\cite{udrescu2020ai}是将深度学习与符号回归结合的代表性工作。该方法利用物理对称性（可加性、可乘性、平移不变性等）作为先验，将高维问题递归分解为低维子问题，再对各子问题执行符号搜索。Sahoo等\cite{sahoo2018learning}提出了EQL（方程学习器），通过在神经网络节点上放置特定的数学算子（如$\sin$、$\cos$、乘法），使网络权重直接编码为公式系数。基于Transformer的符号回归方法\cite{biggio2021neural,kamienny2022end}将方程生成建模为从数据到表达式的序列翻译任务：编码器处理输入-输出数值对，解码器自回归地生成表达式token序列。Kamienny等\cite{kamienny2022end}进一步引入蒙特卡洛树搜索引导解码过程，提升了生成公式的质量。

\textbf{大语言模型驱动的新进展。}SR-LLM\cite{shojaee2025sr}于2025年发表于PNAS，将大语言模型的检索增强生成（RAG）能力引入符号回归。该方法在每轮迭代中从历史搜索结果中检索相似的表达式片段，利用LLM的生成能力提出新的候选公式，显著提升了在复杂方程发现任务上的成功率。

\textbf{基准测试与评价体系。}SRBench\cite{lacava2021srbench}由La Cava等人于2021年建立，涵盖14种符号回归方法和122个基准数据集，是该领域最权威的比较平台。Orzechowski等\cite{orzechowski2018where}的早期基准研究和de Franca等\cite{defranca2023srbench}的SRBench++分别在不同维度上扩展了评价体系。Radwan等\cite{radwan2024comparison}专门比较了近年来新兴算法与传统GP方法的性能差异。Aldeia等\cite{aldeia2024srbench}则呼吁建立下一代符号回归基准，强调需要考察方法在噪声、高维度和分布外泛化等现实条件下的表现。

\textbf{与本课题的关联。}现有符号回归研究——无论是GP方法、神经方法还是LLM方法——主要面向静态回归问题：给定输入-输出对$(x, y)$，发现映射$y = f(x)$。在这一设定中，所有输入变量的地位是对等的，不存在因自回归结构导致的信息竞争。当符号回归被应用于时间序列预测时，自回归滞后特征（如前一时刻的观测值）天然具有极高的预测方差贡献，这在静态回归的基准测试中没有对应的场景。PySR\cite{cranmer2023pysr}在本课题的对齐实验中预测技能分数约为0.076，显著低于KAN教师的0.453，表明直接GP方法在高维时序场景下的搜索效率受限。

\section{Kolmogorov-Arnold网络}

\subsection{数学基础}

Kolmogorov-Arnold表示定理\cite{kolmogorov1957representation,arnold1957functions}是KAN架构的数学基础。1957年，Kolmogorov\cite{kolmogorov1957representation}证明了以下定理：任意定义在有界闭域上的多元连续函数$f(x_1, x_2, \ldots, x_n)$均可精确表示为有限个一元连续函数的嵌套组合与加法：
\begin{equation}
f(x_1, \ldots, x_n) = \sum_{q=0}^{2n} \Phi_q \left( \sum_{p=1}^{n} \phi_{q,p}(x_p) \right)
\end{equation}
其中外层函数$\Phi_q$和内层函数$\phi_{q,p}$均为一元连续函数。Arnold\cite{arnold1957functions}进一步对三变量情形给出了具体的构造证明。该定理的数学意义在于：多元函数的复杂性可以完全归结为一元函数的复杂性和加法运算，不需要乘法或其他多元算子。

\subsection{KAN架构与工作原理}

Liu等\cite{liu2025kan}在2025年ICLR论文中提出的KAN架构直接受Kolmogorov-Arnold表示定理的启发，但做了两个关键的实用化推广。第一，不再要求精确的$(2n+1)$个内层函数，而是允许任意宽度和深度的网络结构；第二，用B样条（B-spline）参数化每条网络边上的一元激活函数，使其可通过梯度下降学习。

KAN与传统MLP的核心区别在于"激活函数放在哪里"。MLP在每个节点上施加固定的激活函数（如ReLU、Sigmoid），节点间的连接承载线性权重。KAN则将可学习的非线性激活函数放在网络边上，每条边是一个由B样条基函数参数化的一元函数，节点仅执行加法汇总。这一设计带来两个后果：首先，KAN的每条边都是一个可独立分析的一元函数映射，具有天然的局部可解释性；其次，当某条边学到的函数接近零函数时，该边可被安全剪除而不影响网络输出——这为结构化稀疏提供了自然的基础。

KAN的参数效率通常优于同等精度的MLP。Liu等\cite{liu2025kan}报告，在多个函数逼近任务上，KAN以百量级参数达到了MLP数千参数的逼近精度。Kratsios和Furuya\cite{kratsios2024kan}从理论上分析了KAN的逼近和学习保证，证明其在函数导数逼近方面也具有优势。Howard等\cite{howard2024finite}提出了有限基KAN变体，通过域分解提升了其在物理信息问题上的效率。

\subsection{符号提取流程}

KAN符号提取是一个从数值B样条到解析表达式的多步转换过程，也是KAN可解释性优势的具体实现路径\cite{liu2025kan,liu2025kan2}。该流程包含以下步骤：

\textbf{第一步：常规训练（warmup）。}使用标准的均方误差损失训练KAN至收敛，此时所有边上的B样条函数都处于自由拟合状态。

\textbf{第二步：稀疏化（sparsification）。}在损失函数中引入L1正则化和熵正则化项。L1正则化惩罚每条边的B样条系数绝对值之和，驱动不重要的边趋向零函数。熵正则化惩罚边的激活分布的均匀性——直观地，如果一条边在所有输入上的激活幅值都很小，其信息贡献为零，应被剪除。两种正则化的联合作用使KAN在训练过程中自动形成稀疏结构。

\textbf{第三步：结构化剪枝（pruning）。}基于阈值搜索，将激活幅值低于阈值的边和孤立节点从网络中移除。剪枝后的KAN是原始网络的一个稀疏子图，仅保留信息贡献显著的边。

\textbf{第四步：逐边符号拟合（symbolic fitting）。}对每条存活的边，用一组候选符号函数（如$x, x^2, x^3, \sin(x), \cos(x), |x|$等）逐一拟合其B样条曲线。选择拟合$R^2$最高且超过预设阈值的候选函数替代原始B样条。

\textbf{第五步：公式合成（formula composition）。}将所有已符号化的边按照网络拓扑组合，得到从输入到输出的完整解析表达式。使用SymPy进行化简，生成最终的闭式公式。

\subsection{KAN 2.0与科学应用}

Liu和Tegmark\cite{liu2025kan2}在2025年Physical Review X上的工作（KAN 2.0）系统展示了KAN在科学方程发现中的能力。该工作引入了MultKAN（支持乘法节点的KAN变体）和kanpiler（将KAN编译为可执行函数的工具），并在12个经典物理方程上成功执行了符号提取。关键案例包括：从数值数据中恢复Kolmogorov湍流定标律（能谱的五三次幂律）、Stokes定律、以及特殊函数的级数展开。这些案例表明KAN符号提取在"干净"物理问题上具有可靠的方程发现能力。

然而，上述所有成功案例都具有一个共同的前提条件：输入变量之间不存在信息层级差异。每个物理量（如流体速度、粘度、球体半径）对输出都有不可替代的独立贡献。在时间序列预测场景中，自回归滞后特征通常解释了90\%以上的短期方差，与外生物理变量之间形成了极不对称的信息竞争关系。这一结构性差异是KAN符号提取从物理方程发现迁移到时间序列预测时的核心挑战，也是本课题研究的出发点。

\section{MCKAN与风光预测参考}

Chen等\cite{chen2025mckan}提出的MCKAN（Multi-scale Convolutional Kolmogorov-Arnold Network）是KAN在风光功率多步预测中的代表性应用，发表于Engineering Applications of Artificial Intelligence。MCKAN的架构设计和方法论选择对本课题具有直接的参考价值。

\textbf{架构设计。}MCKAN采用多尺度卷积模块提取不同时间粒度的特征模式，通过EAA（Enhanced Additive Attention）注意力机制融合多尺度特征，并以KAN层作为最终的预测映射。该设计的核心思想是利用卷积的局部感受野和KAN的函数逼近能力互补——卷积擅长提取局部时间模式（如日内周期），KAN擅长学习全局非线性映射。MCKAN还引入了CQALA（改进的Lemming优化器+自适应注意力融合）进行超参数优化和注意力权重调整。

\textbf{方法论选择。}MCKAN在数据预处理上采用Pearson相关系数进行特征选择，保留与预测目标相关性超过阈值的特征。归一化采用Min-Max缩放，将特征值映射到$[0, 1]$区间。预测目标覆盖15分钟至45分钟的多步预测（multi-step forecasting）。消融实验验证了各模块（多尺度卷积、EAA注意力、CQALA优化）的独立贡献。

\textbf{本课题的对应关系。}本课题的方法设计在多个层面与MCKAN形成参照：（1）MCKAN的多步预测设计对应本课题的多步差分目标（delta targets at various horizons）；（2）MCKAN的Pearson相关特征选择对应本课题的VER/FAR可辨识性诊断——两者都在回答"哪些特征对预测有用"，但VER/FAR进一步追问"有用的特征能否进入显式公式"；（3）MCKAN的风光子任务结构对应本课题的S3结构化分解；（4）MCKAN的模块消融对应本课题的S2自回归阻断消融——两者都通过移除特定组件观察系统行为变化，但S2阻断实验的目标不是验证模块贡献，而是检验信息竞争机制。

\textbf{与本课题的差异。}MCKAN聚焦于通过架构创新提升预测精度，其评价体系以RMSE、MAE、$R^2$等精度指标为主。本课题则聚焦于从训练后的KAN中提取符号公式这一下游任务，评价体系除精度指标外还包含VER、FAR、活跃边数等可辨识性指标。两项工作在KAN的使用方式上也存在本质差异：MCKAN将KAN作为端到端预测模型的组成部分，不执行后验符号提取；本课题将KAN作为"先学习后提取"管线的教师模型，符号提取是方法的核心步骤而非可选的事后分析。

\section{物理信息学习与可解释性}

将物理先验知识嵌入机器学习模型，是提升模型可解释性和物理一致性的另一条技术路径。

\textbf{物理信息神经网络（PINN）。}Raissi等提出的PINN框架将偏微分方程的残差作为损失函数的正则化项，使神经网络在拟合数据的同时满足已知的物理约束。Misyris等\cite{misyris2020physics}将PINN引入电力系统动态仿真，用于求解摇摆方程（swing equation）。Stiasny等\cite{stiasny2021learning}进一步提出了无数据训练的PINN方案——仅依赖物理方程作为监督信号训练网络，消除了对大规模仿真数据的依赖。PINNSim\cite{pinnsim2023}将PINN方法工程化为可用的电力系统动态仿真器。

\textbf{PIKAN。}Shuai和Li\cite{shuai2024pikan}将物理信息学习与KAN架构结合，提出PIKAN（Physics-Informed KAN），用于电力系统动力学模拟。PIKAN利用KAN边上B样条的灵活性编码物理约束，在摇摆方程求解中达到了优于PINN的精度。这一工作表明KAN架构与物理先验知识的结合是可行的，但PIKAN的目标是求解已知方程的数值解，而非从数据中发现未知方程。

\textbf{可解释人工智能（XAI）在能源领域的应用。}Machlev等\cite{machlev2022explainable}在2022年发表于Energy and AI的综述中系统梳理了XAI技术在能源和电力系统中的应用现状。该综述将可解释性方法按层次分为三类：（1）特征重要性方法（如SHAP、LIME），提供输入变量对输出的边际贡献排序；（2）注意力可视化方法，通过注意力权重矩阵展示模型在时间和特征维度上的关注模式；（3）符号回归方法，直接输出闭式数学公式。在这三个层次中，特征重要性和注意力可视化属于"近似事后解释"——它们描述模型行为的统计特征，但不等于模型本身；符号回归则属于"内建可解释性"——公式本身就是模型，不存在解释与模型之间的忠实度（faithfulness）问题。

\textbf{KAN在可解释能源预测中的应用。}Panczyk等\cite{panczyk2024kan}在核工程领域验证了KAN-to-symbol的转换路径，从核反应堆热工水力数据中成功提取了包含物理变量的解析公式。Mubarak等\cite{mubarak2024kan}将采用拟牛顿优化的KAN应用于风电场功率预测，在预测精度上取得了与深度学习方法可比的结果，但未进一步执行符号提取步骤。

\textbf{可解释性的层次分析。}综合上述工作，可以将当前的可解释性方法按"解释强度"从弱到强排列：统计归因（SHAP/LIME）$\to$ 注意力可视化 $\to$ 决策规则（决策树）$\to$ 符号公式。前三种方法提供的是关于模型行为的间接信息，其解释的忠实度取决于近似方法的精度。符号公式提供的是模型本身的直接表达——公式即模型，不存在忠实度损失。KAN的符号提取能力使其在理论上能够达到最高层次的可解释性。然而，正如本课题所揭示的，从高精度KAN教师到符号公式的转换在实践中会遭遇信息损失——尤其是当自回归滞后特征的存在导致物理变量在稀疏化中被剪除时，最终公式虽然在形式上是"符号可解释的"，但其物理含义已被掏空。这一精度-可解释性-物理忠实度的三角矛盾是本课题研究的核心问题。

\section{本章小结}

本章从五个方向梳理了与本课题相关的研究进展。负荷与净负荷预测领域已形成从统计方法到深度学习的成熟技术栈，但高精度方法均为黑箱。符号回归领域从GP到神经方法到LLM驱动方法不断演进，但研究对象集中于静态回归问题，尚未涉及时间序列中的自回归竞争机制。KAN架构及其符号提取能力已在物理方程发现中得到验证，但成功案例均不包含自回归特征。MCKAN等工作将KAN应用于风光预测并取得了精度上的进展，但未执行符号提取。物理信息学习和XAI领域提供了丰富的可解释性技术手段，但尚未解决"高精度模型的符号公式中物理变量为何消失"这一问题。

上述五条脉络共同指向一个未被填补的研究空白：在KAN符号提取用于时间序列预测的场景中，没有工作从机制层面研究物理变量的符号可辨识性问题。具体而言，三个问题尚无答案：（1）自回归滞后特征的捷径竞争是否是物理变量被剪除的主要机制？（2）这一机制能否通过可控的干预实验（而非仅凭观察性证据）得到验证？（3）是否存在结构化的方法恢复物理变量的可辨识性？本课题正是围绕这三个问题展开研究。
